\documentclass[pdflatex,sn-mathphys-num]{sn-jnl}% Math & Physical Sciences, numbered

%%%% Standard packages
\usepackage{graphicx}%
\usepackage{multirow}%
\usepackage{amsmath,amssymb,amsfonts}%
\usepackage{amsthm}%
\usepackage{mathrsfs}%
\usepackage[title]{appendix}%
\usepackage{xcolor}%
\usepackage{textcomp}%
\usepackage{manyfoot}%
\usepackage{booktabs}%
\usepackage{algorithm}%
\usepackage{algorithmicx}%
\usepackage{algpseudocode}%
\usepackage{listings}%
\usepackage{siunitx}%
%% Custom units for geodesy/gravimetry: gal, nanostrain, and year.
\DeclareSIUnit{\gal}{Gal}%
\DeclareSIUnit{\strain}{strain}%
\DeclareSIUnit{\year}{yr}%

\graphicspath{{figures/}}

%% Code-listing style (Methods listings + Appendix listings).
\lstdefinestyle{py}{language=Python,basicstyle=\small\ttfamily,
  keywordstyle=\color{blue!60!black},commentstyle=\color{gray},
  stringstyle=\color{green!45!black},showstringspaces=false,
  numbers=left,numberstyle=\tiny\color{gray},frame=single,breaklines=true}

\theoremstyle{thmstyleone}%
\newtheorem{theorem}{Theorem}%
\newtheorem{proposition}[theorem]{Proposition}%
\theoremstyle{thmstyletwo}%
\newtheorem{example}{Example}%
\newtheorem{remark}{Remark}%
\theoremstyle{thmstylethree}%
\newtheorem{definition}{Definition}%

\raggedbottom

\begin{document}

%% ---------------------------------------------------------------------
%% TITLE.  Working title for Journal of Geodesy (geodesy/geodynamics focus).
%% Options:
%%  (1) Geodetic strain rates, gravity-field gradients and seismicity of the
%%      Western Anatolian graben system mapped by machine learning and PyGMT
%%  (2) Machine-learning fusion of GNSS strain rates and gravity-field data for
%%      geodynamic characterization of the Aegean extensional province, Turkiye
%%  (3) Crustal extension of the Western Anatolian grabens from geodetic velocity,
%%      strain-rate and geoid fields: a machine-learning and PyGMT analysis
%% ---------------------------------------------------------------------
%\title[ML geodesy of the Western Anatolian grabens]{[REPLACE with chosen title --- working: \emph{Geodetic strain rates, gravity-field gradients and seismicity of the Western Anatolian graben system mapped by machine learning and PyGMT}]}

\title[MLfusion of geodetic and gravity data in western Anatolia]{\emph{Machine-learning fusion of geodetic and gravity data for modelling crustal strain and seismicity in Western Anatolian graben
system}}

%%=============================================================%%
\author*[1]{\fnm{Polina} \sur{Lemenkova}}
\email{lemenkova@itu.edu.tr}
\author[1]{\fnm{Abdullah Can} \sur{Zülfikar}}
\email{aczulfikar@itu.edu.tr}

\affil*[1]{\orgdiv{Disaster and Emergency Management Department, Institute of Earthquake Engineering and Disaster Management}, \orgname{Istanbul Technical University}, \orgaddress{\street{İTÜ Ayazağa, Maslak Kampüsü}, \city{İstanbul}, \postcode{34469}, \country{Türkiye}}}

\abstract{[REPLACE: field intro --- geodetic monitoring of active crustal
deformation and its link to seismic hazard in extensional tectonic provinces.]
[REPLACE: background --- the Western Anatolian Extensional Province (Aegean graben
system) is one of the most rapidly extending continental regions on Earth, with
open geodetic, gravity and seismological archives now available.] [REPLACE: problem
/ gap --- its geodetic deformation, gravity-field signature and seismicity have not
been jointly analysed and mapped with a unified, open, data-driven workflow.] Here
we [REPLACE: what this paper does --- use machine learning to relate GNSS/InSAR
velocity and strain-rate fields, gravity/geoid gradients and seismicity across the
grabens, and render the results as continuous maps with GMT/PyGMT]. [REPLACE: key
quantified result vs prior knowledge --- model skill {\normalfont\unboldmath
$R^2=[REPLACE]$}, strain-rate magnitudes and the geodetic--gravity--seismicity
correspondence, with numbers.] [REPLACE: broader context --- a reproducible,
open-data template for geodetic--geodynamic characterization that transfers to other
extensional provinces.]}

\keywords{geodesy, crustal deformation, strain rate, gravity field, machine learning, Western Anatolia, Aegean graben system, PyGMT}

\maketitle

% =====================================================================
\section{Introduction}\label{sec1}
Space geodesy has transformed the study of active continental deformation, turning
the slow strain of the lithosphere into a directly measurable field. Global
Navigation Satellite System (GNSS) networks and interferometric synthetic aperture
radar (InSAR) now resolve surface velocities at the millimetre-per-year level over
entire orogens, while satellite gravimetry and high-resolution gravity-field models
constrain the mass distribution and crustal structure that accompany that
deformation \cite{Wright2004,MassonnetFeigl1998,Burgmann2000,Elliott2016,Tapley2004,Pavlis2012}. Quantifying how interseismic
strain accumulates, where it localises, and how it relates to the release of
seismic energy is a central problem of geodetic science, because the same velocity
and gravity fields that describe the present-day kinematics also condition the
seismic hazard faced by the populations living above active faults
\cite{Hussain2018,Reilinger2006,Lemenkova202316,Weiss2020,Ergintav2023}. Extensional provinces, where the
crust is actively thinning across networks of normal faults, are an especially
demanding test of this capability, since their deformation is distributed,
predominantly north--south, and entangled with strong non-tectonic signals.

The eastern Mediterranean is one of the most thoroughly instrumented natural
laboratories for this purpose. Decades of GNSS campaigns have mapped the westward
escape of the Anatolian microplate between the North and East Anatolian fault zones
and its south-westward extrusion towards the Hellenic trench, yielding a
present-day velocity field referenced to stable Eurasia
\cite{McClusky2000,Reilinger2006,Nocquet2012,Kurt2023}. Building on these networks,
combined GNSS and Sentinel-1 InSAR analyses have produced continuous,
high-resolution velocity and strain-rate fields for the whole of Anatolia, isolating
the localised strain along the major strike-slip faults from the more diffuse
deformation of the Aegean domain \cite{England2016,Weiss2020,Kreemer2003,Hollenstein2008,Berardino2002,Ferretti2001,Hussain2018}. In
parallel, the gravity field and geoid derived from satellite missions and global
models such as EGM2008 have been used to image lithospheric structure, sediment
thickness and isostatic state in tectonically comparable settings, complementing the
purely kinematic picture with information on subsurface mass
\cite{Pavlis2012,Sandwell2014,Lemenkova202216,Foerste2014,Ince2019,Hirt2013,HofmannWellenhof2006,Tapley2004,Farr2007}. At the same time,
data-driven and machine-learning methods have moved from peripheral tools to
mainstream instruments of the solid-Earth sciences, learning the non-linear
relationships among heterogeneous geophysical observables that fixed parametric
forms approximate only crudely
\cite{Bergen2019,Reichstein2019,MousaviBeroza2023,Lemenkova202518,LeCun2015,DeVries2018,Dramsch2020,Lary2016,Karpatne2019,Breiman2001,ChenGuestrin2016}.

The Western Anatolian Extensional Province, the Aegean graben system, is the third
major seismic province of T\"{u}rkiye and the locus of the region's most rapid
continental extension. Driven by the roll-back of the Hellenic subduction zone and
the westward motion of Anatolia, the province is opening in a roughly north--south
direction at rates of the order of 20~mm/yr, accommodated by predominantly
east--west-trending normal faults that bound a family of elongated grabens --- the
Gediz (Ala\c{s}ehir), B\"{u}y\"{u}k Menderes, K\"{u}\c{c}\"{u}k Menderes and Simav
basins among them
\cite{McKenzie1972,Jackson1994,Aktug2009,LePichonAngelier1979,SengorYilmaz1981,Bozkurt2001,Jolivet2013,BozkurtSozbilir2004,TenVeen2009}. Towards
the eastern Aegean coast the normal-fault network gives way to the transtensional
\.{I}zmir--Bal\i kesir transfer zone, where north-east-trending faults carry a
significant right-lateral component
\cite{Uzel2012,Taymaz2007,Lemenkova202215,Nissen2022,Kocyigit2005}. The province has repeatedly produced
damaging normal-faulting earthquakes of magnitude up to about seven, including the
1969 Ala\c{s}ehir event and, more recently, moderate shocks such as the 2019
Acipayam and Acig\"{o}l ruptures and the 2020 Samos sequence, several of which left
clear coseismic signatures in Sentinel-1 interferograms
\cite{EyidoganJackson1985,Yang2020,Nissen2022,Storchak2013}. Recent InSAR
time-series studies have shown that strain in these grabens is not confined to the
main bounding faults but is distributed across numerous smaller and antithetic
structures, favouring a continuum rather than a rigid-block description of the
deformation \cite{Diercks2024,Weiss2020,Aktug2009}.

Despite this wealth of observation, the geodetic, gravimetric and seismological
descriptions of the Aegean graben system have largely been developed in isolation.
GNSS studies resolve the regional velocity and strain-rate field but are too sparse
to constrain deformation at the scale of individual grabens; InSAR delivers high
spatial resolution but is weakly sensitive to the dominant north--south motion and
is swamped, in the basins, by anthropogenic subsidence from groundwater extraction
and by topographic and atmospheric artefacts
\cite{Wright2004,Diercks2024,Dogan2022,Imamoglu2022}. Gravity-field information, where used,
is rarely co-registered with the strain-rate and seismicity fields, and the
spatially continuous relationship between crustal extension, its gravimetric
expression and the distribution of earthquakes has not been established with a
single, open and reproducible workflow. No study has yet brought the geodetic
velocity and strain-rate fields, the gravity and geoid signature, and the
instrumental seismicity of the Western Anatolian grabens into one data-driven
framework that predicts and maps them as co-registered, continuous fields
\cite{Diercks2024,Weiss2020,Ergintav2023}.

Closing this gap matters well beyond methodology. A joint, spatially continuous view
of strain rate, gravity and seismicity would let the rapidly evolving deformation of
a densely populated province be monitored in a reference-frame-consistent way, would
sharpen the geodynamic interpretation of how extension is partitioned between
localised and distributed faulting, and would provide the spatially resolved strain
information on which seismic-hazard assessment in western T\"{u}rkiye ultimately
depends \cite{Hussain2018,Ergintav2023,Lemenkova202606,Briole2021,Kurt2023}. Because the required
records --- open GNSS velocity solutions, freely processed Sentinel-1 InSAR
products, global gravity-field models, instrumental earthquake catalogues and
near-global terrain data --- already exist for many extensional settings, a workflow
validated on the Aegean grabens would transfer directly to other actively extending
regions \cite{Kurt2023,Lazecky2020,Pavlis2012}.

The objective of this study is therefore to predict and map the geodetic
strain-rate field, the associated gravity and geoid signature, and their relation to
seismicity along the Western Anatolian Extensional Province from open-access data
using machine learning, to render the results as continuous, co-registered maps with
GMT and PyGMT, and to quantify the associated uncertainty. The analysis relies only
on openly available datasets and on open-source Python and GMT/PyGMT tooling, so that
it is fully reproducible without any field campaign or restricted data
\cite{Wessel2019,Harris2020,Pedregosa2011}. The main contributions are the
following:
\begin{itemize}
  \item a unified, open, fully data-driven workflow that joins machine-learning
        regression of geodetic and gravimetric fields in Python to cartographic
        rendering in GMT/PyGMT, applied to a continental extensional province;
  \item continuous, co-registered maps of the GNSS/InSAR velocity field, the
        strain-rate field and the gravity/geoid gradients across the Gediz,
        B\"{u}y\"{u}k Menderes and neighbouring grabens;
  \item a quantified relationship between crustal extension, its gravimetric
        expression and instrumental seismicity, with an uncertainty treatment under
        spatially grouped cross-validation that locates where the prediction
        uncertainty arises;
  \item a reproducible, open-data template that transfers directly to other
        instrumented extensional settings.
\end{itemize}
The remainder of the paper is organised as follows. Section~\ref{sec:study}
describes the study area and the open-access datasets; Section~\ref{sec:methods}
sets out the geodetic and gravimetric background and the modelling, mapping and
validation procedures; Section~\ref{sec:results} presents the results;
and Sections~\ref{sec:discussion} and~\ref{sec:conclusions} interpret them and
conclude.

% =====================================================================
\section{Study area and data}\label{sec:study}
This section establishes the geological and observational basis of the analysis. It
first describes the seismotectonic and geodynamic setting of the Western Anatolian
Extensional Province, then summarises its instrumental seismicity and focal
mechanisms, and finally enumerates the open-access datasets from which the
predictors and targets of the machine-learning workflow are built. All datasets are
freely available, so that the analysis is fully reproducible without any field
campaign or restricted data.

\subsection{Seismotectonic and geodynamic setting}\label{sec:setting}
The study region is the Western Anatolian Extensional Province (WAEP), the system of
active grabens that occupies the western third of T\"{u}rkiye between roughly
$26^{\circ}$ and $31^{\circ}$E and $36^{\circ}$ and $40.5^{\circ}$N. Its deformation
is driven by two coupled processes: the westward escape of the Anatolian microplate
between the North and East Anatolian fault zones, and the south-westward roll-back of
the Hellenic subduction front, which together place the overriding Aegean lithosphere
in tension \cite{McKenzie1972,Reilinger2006,Jackson1994,LePichonAngelier1979,Jolivet2013,Taymaz2007}. The combined
effect is rapid, broadly north--south extension at rates of the order of 20~mm/yr,
accommodated by a dense network of predominantly east--west-trending normal faults
that bound a family of elongated sedimentary basins
\cite{Aktug2009,McClusky2000,SengorYilmaz1981,Bozkurt2001,TenVeen2009}. The most prominent of these are the Gediz
(Ala\c{s}ehir) and B\"{u}y\"{u}k Menderes grabens in the centre of the province, the
K\"{u}\c{c}\"{u}k Menderes graben between them, and the Simav graben to the north,
together with smaller basins such as the Denizli, \c{C}ivril and Acig\"{o}l
depressions to the east \cite{BozkurtSozbilir2004,Kocyigit2005,Emre2018,McKenzie1978,Diercks2024}.

The internal architecture of these grabens records a protracted extensional history.
In the Gediz graben, a low-angle detachment fault exhumed the Menderes Massif during
the Miocene and is generally regarded as inactive, while a younger, high-angle Gediz
graben-bounding fault subsequently formed in its hangingwall and now controls the
southern basin margin \cite{BozkurtSozbilir2004,Kent2016,Diercks2024}. Geodetic and
geomorphic studies indicate that extension is not localised on these principal
bounding faults alone but is distributed across numerous smaller, antithetic and
secondary structures within the basins, a pattern more consistent with a continuum
than with a rigid-block description of the deformation
\cite{Aktug2009,Weiss2020,Diercks2024}. Towards the eastern Aegean coast the
normal-faulting regime passes into the transtensional \.{I}zmir--Bal\i kesir transfer
zone, where north-east-trending faults carry a significant right-lateral component
and reorganise the strain field \cite{Uzel2012,Nissen2022,Lemenkova202215,Taymaz2007,Ergintav2023}.
The regional setting --- shaded relief, the principal grabens and active-fault
traces, the plate-boundary and transfer zones, the open geodetic sites and the
instrumental seismicity --- is summarised in Fig.~\ref{fig:studyarea}.

\begin{figure}[htbp]\centering
	\includegraphics[width=0.95\textwidth]{fig01_studyarea}
  \caption{Study-area and geodynamic setting of the Western Anatolian Extensional
  Province (Aegean graben system): shaded relief, principal grabens (Gediz,
  B\"{u}y\"{u}k and K\"{u}\c{c}\"{u}k Menderes, Simav), active normal faults,
  plate-boundary and \.{I}zmir--Bal\i kesir transfer zones, open GNSS sites and
  instrumental seismicity, with a locator inset. Software used for plotting figure:
  GMT/PyGMT \texttt{<ver>}. Data: SRTM/GEBCO relief; active-fault traces; GNSS
  network; earthquake catalogue. Source: authors.}\label{fig:studyarea}
\end{figure}

\subsection{Regional seismicity and focal mechanisms}\label{sec:seis}
The instrumental seismicity of the province reflects this distributed extensional
regime. Earthquakes concentrate along the graben-bounding faults and cluster in the
upper crust, consistent with brittle normal faulting in actively thinning continental
lithosphere \cite{EyidoganJackson1985,Jackson1994,Storchak2013}. The province has
repeatedly produced damaging earthquakes of magnitude up to about seven: the 1969
Ala\c{s}ehir (Sar\i g\"{o}l) event on the eastern Gediz graben remains the reference
large normal-faulting rupture, and the more recent instrumental record includes the
2019 Acipayam and Acig\"{o}l earthquakes on blind and antithetic basin faults and the
2020 Samos sequence offshore \.{I}zmir
\cite{EyidoganJackson1985,Yang2020,Nissen2022,Storchak2013}. Several of these moderate
events left coseismic displacement signatures detectable in Sentinel-1 interferograms,
confirming that the mapped faults are active and that seismic release is spatially
correlated with the interseismic deformation field \cite{Diercks2024,Lazecky2020,Yang2020}.

The kinematics of these earthquakes are dominated by normal faulting on
east--west-striking planes, with subordinate strike-slip mechanisms along the
north-east-trending faults of the \.{I}zmir--Bal\i kesir transfer zone
\cite{Uzel2012,Nissen2022,Lemenkova202205,Taymaz2007,Jackson1994}. Centroid-moment-tensor
solutions provide a consistent, openly available description of these source
geometries and, rendered as focal-mechanism beachballs over shaded relief, expose the
along-strike variation in faulting style across the province
\cite{Dziewonski1981,Ekstrom2012,Lemenkova202404,Taymaz2007,Uzel2012,Ergintav2023}. The
regional instrumental seismicity and the focal mechanisms of the principal events are
summarised in Fig.~\ref{fig:beachballs}, which anchors the data-driven analysis that
follows in the observed pattern of seismic release.

\begin{figure}[htbp]\centering
	\includegraphics[width=0.95\textwidth]{fig02_seismicity_beachballs}
  \caption{Regional instrumental seismicity scaled by magnitude and coloured by focal
  depth, with GCMT focal-mechanism beachballs for the principal normal-faulting events
  (e.g.\ the 2020 Samos and 2019 Acig\"{o}l earthquakes), over grey-scale shaded
  relief. Software used for plotting figure: GMT/PyGMT \texttt{<ver>}. Data:
  instrumental earthquake catalogue; GCMT moment tensors; SRTM/GEBCO relief.
  Source: authors.}\label{fig:beachballs}
\end{figure}

\subsection{Open-access datasets}\label{sec:data}
Five open-access datasets underpin the analysis (Table~\ref{tab:data}). Surface
kinematics are taken from published GNSS velocity solutions for T\"{u}rkiye and the
Aegean, referenced to a stable Eurasia frame, which provide the long-wavelength
north--south and east--west velocity components that constrain the regional strain
field \cite{Kurt2023,England2016,Nocquet2012}. These are complemented by
freely processed Sentinel-1 InSAR products from the automated LiCSAR/LiCSBAS
processing chain, which deliver high-spatial-resolution line-of-sight velocity fields
and their decomposition into vertical and east--west components across the grabens
\cite{Lazecky2020,Morishita2020,Berardino2002,Ferretti2001,Goldstein1988,Rosen2000,Weiss2020}. The gravity field and geoid are obtained
from a high-degree global gravity-field model, providing the gravity disturbance and
geoid-gradient fields that carry information on crustal and lithospheric structure
beneath the basins \cite{Pavlis2012,Foerste2014,Lemenkova202216,Ince2019,Hirt2013,Tapley2004,Sandwell2014}.
Event hypocentres, magnitudes and focal depths are drawn from open instrumental
earthquake catalogues together with centroid-moment-tensor solutions, and topography
and bathymetry are taken from the SRTM and GEBCO global grids, complemented where
needed by the ALOS global DEM, and processed into terrain models with GMT/PyGMT and
the scientific Python stack
\cite{Storchak2013,Ekstrom2012,Farr2007,GEBCO2023,Tadono2014,Wessel2019}.

The same script-driven cartographic approach has been applied to the topography,
gravity field and seismicity of tectonically comparable settings, including
fold-and-thrust belts, subduction trenches and continental rifts, and the open-source
machine-learning and image-analysis libraries adopted here have likewise been
exercised in comparable Earth-observation pipelines, underpinning the reproducibility
of the workflow
\cite{Wessel2019,Sandwell2014,Lemenkova202316,Harris2020,Pedregosa2011,Lemenkova202301,Farr2007,Virtanen2020}.
Built-up areas, which generate spurious deformation signals in the InSAR fields, are
identified and masked using the openly available Global Human Settlement Layer
\cite{Pesaresi2023,Imamoglu2022,Diercks2024}. The distribution of the resulting sample
in geodetic, gravimetric and seismological space, together with its principal coverage
gaps, is shown in Fig.~\ref{fig:datadist}.

\begin{table}[htbp]
  \caption{Open-access datasets used in this study, with provider, version and access
  (all open access).}\label{tab:data}
  \begin{tabular*}{\textwidth}{@{\extracolsep\fill}lll}
    \toprule
    Dataset & Provider / version & Access \\
    \midrule
    Geodetic velocities (GNSS) & Regional GNSS solutions, stable-Eurasia frame & open access \\
    Velocity / strain (InSAR)  & Sentinel-1 LiCSAR / LiCSBAS products & open access \\
    Gravity field / geoid      & Global gravity-field model (high degree) & open access \\
    Earthquake catalogue       & Instrumental catalogue + GCMT tensors & open access \\
    Topography / bathymetry    & SRTM / GEBCO grids (with ALOS DEM) & open access \\
    \botrule
  \end{tabular*}
\end{table}

\begin{figure}[htbp]\centering
  \includegraphics[width=0.95\textwidth]{fig03_data_distribution}
  \caption{Distribution of the geodetic, gravimetric and seismological sample used for
  training and evaluation: GNSS and InSAR coverage, velocity-magnitude and
  depth--magnitude histograms, and spatial sampling across the province. Software used
  for plotting figure: Python~3.12, Matplotlib~3.10. Data: as in
  Table~\ref{tab:data}. Source: authors.}\label{fig:datadist}
\end{figure}

% =====================================================================
\section{Methods}\label{sec:methods}
This study predicts the geodetic strain-rate field, the associated gravity and geoid
signature, and their relation to seismicity along the Western Anatolian Extensional
Province, and maps the results through a single, open and reproducible pipeline that
joins Python machine learning to GMT/PyGMT cartography. The subsections below follow
the workflow in turn: the overall framework, the geodetic and gravimetric relations
that motivate the predictor features, the preprocessing and feature construction, the
model and its training, the experimental configuration, the cartographic rendering,
and the validation and uncertainty analysis. Throughout, the standing project
defaults hold: only open-access data are used, the tools are Python with its
scientific and machine-learning stack together with GMT/PyGMT, and the methods are
modelling, mapping and data analysis.

\subsection{Methodological framework}\label{sec:framework}
The workflow proceeds from the open-access datasets of Section~\ref{sec:study} through
quality control and feature engineering, model training and cross-validated tuning,
prediction on a regular geographic grid, and finally field mapping and validation. The
GNSS and InSAR velocity solutions, the gravity-field model, the earthquake catalogue
and the terrain grids enter on the input side; a gradient-boosted regression ensemble
learns the mapping from source, position, terrain and gravimetric predictors to the
geodetic and gravimetric targets; and the trained model is evaluated on a regular grid
and rendered as continuous, co-registered maps. The complete sequence is summarised
schematically in Fig.~\ref{fig:workflow} and as pseudocode in
Algorithm~\ref{alg:pipeline}, with the dashed feedback loop in the figure denoting the
cross-validated hyper-parameter tuning \cite{Pedregosa2011,ChenGuestrin2016,Reichstein2019}.

\begin{figure}[htbp]\centering
  \includegraphics[width=0.95\textwidth]{fig04_workflow}
  \caption{Methodological workflow from open-access geodetic, gravimetric and seismic
  data, through Python feature engineering and machine-learning modelling, to
  GMT/PyGMT field mapping and validation; the dashed loop denotes the cross-validated
  hyper-parameter tuning. Software used for plotting figure: Python~3.12, Matplotlib~3.10.
  Data: schematic (no data plotted). Source: authors.}\label{fig:workflow}
\end{figure}

Figure~\ref{fig:workflow} traces this pipeline from left to right: the open-access
inputs enter on the left, the central block performs quality control, feature
engineering and model training with cross-validated tuning, and the right-hand block
carries the trained model through grid prediction and GMT/PyGMT rendering to the
validated field maps. The same procedure is stated as explicit pseudocode in
Algorithm~\ref{alg:pipeline}.

\begin{algorithm}
\caption{End-to-end geodetic-field prediction and mapping pipeline.}\label{alg:pipeline}
\begin{algorithmic}[1]
\Require GNSS velocities $V$, InSAR velocity/strain fields $I$, gravity model $G$, earthquake catalogue $C$, terrain grids $T$
\Ensure Trained model $f$, predicted geodetic/gravimetric fields, co-registered maps
\State Merge $V$, $I$, $G$, $C$ and $T$ onto a common geographic grid
\State Quality-control records; mask built-up areas and low-coherence pixels
\State Form predictors $\mathbf{x}=(\lambda,\varphi,v_e,v_n,\text{depth},\log T,\Delta g)$
\State Set targets $y$ (strain-rate scalar, gravity gradient, seismicity proxy)
\State Split into training and held-out test sets, grouped by spatial block
\State Standardise predictors using training-set statistics only
\State Tune hyper-parameters by spatially grouped $k$-fold cross-validation
\State Fit the gradient-boosted ensemble $f$ on the training set
\State Evaluate $R^2$, RMSE and MAE on the held-out test set
\State Predict the target fields on a regular $0.02^{\circ}$ grid
\State Render predicted-field and co-registered maps with GMT/PyGMT
\end{algorithmic}
\end{algorithm}

\subsection{Geodetic and gravimetric background}\label{sec:geobg}
The quantities predicted in this study, and the relations that motivate the predictor
features, rest on a small set of well-established geodetic and gravimetric
expressions, summarised here to fix notation. Active crustal deformation is described
by the horizontal velocity field $\mathbf{v}=(v_e,v_n)$ measured at the surface by
GNSS and InSAR. Its spatial gradients define the strain-rate tensor
\begin{equation}
\dot{\varepsilon}_{ij}=\tfrac{1}{2}\!\left(\frac{\partial v_i}{\partial x_j}+\frac{\partial v_j}{\partial x_i}\right),
\qquad i,j\in\{e,n\},
\label{eq:strain}
\end{equation}
where $v_e$ and $v_n$ are the east and north velocity components
(\si{\milli\metre\per\year}) and $x_e,x_n$ are the corresponding map coordinates, so
that each $\dot{\varepsilon}_{ij}$ has units of \si{\nano\strain\per\year}
\cite{England2016,Weiss2020,Kreemer2003,Hollenstein2008}. Because the individual tensor components depend on the
chosen coordinate frame, the deformation is summarised by frame-invariant scalars. The
dilatation rate $\dot{\varepsilon}_{v}=\dot{\varepsilon}_{ee}+\dot{\varepsilon}_{nn}$
measures area change --- positive for the extension that opens the grabens --- while
the second invariant of the strain-rate tensor,
\begin{equation}
\dot{\varepsilon}_{\mathrm{II}}=\sqrt{\dot{\varepsilon}_{ee}^{2}+\dot{\varepsilon}_{nn}^{2}+2\,\dot{\varepsilon}_{en}^{2}},
\label{eq:invariant}
\end{equation}
provides a single non-negative measure of the total strain-rate magnitude that is used
here as the principal geodetic target \cite{Hussain2018,Kurt2023,Ergintav2023}.

The gravity field complements this kinematic picture with information on subsurface
mass. Relative to a reference ellipsoid, the gravity disturbance is
\begin{equation}
\delta g(\mathbf{r})=g(\mathbf{r})-\gamma(\mathbf{r}),
\label{eq:gravity}
\end{equation}
where $g(\mathbf{r})$ is the observed gravity and $\gamma(\mathbf{r})$ the normal
gravity of the reference field at position $\mathbf{r}$, both in
\si{\milli\gal}; the horizontal gradients of $\delta g$ and of the geoid undulation
track lateral density contrasts such as the sediment-filled basins and the exhumed
metamorphic massifs that flank them \cite{Pavlis2012,Foerste2014,Sandwell2014}. The
gravity disturbance, its horizontal gradient and the geoid gradient enter the workflow
both as gravimetric predictors and, in the case of the gravity gradient, as a mapped
target field \cite{Pavlis2012,Foerste2014,Lemenkova202216,Sandwell2014,Tapley2004,Wessel2019}.

\subsection{Data preprocessing and feature engineering}\label{sec:prep}
The predictors and targets are assembled from the merged GNSS, InSAR, gravity,
catalogue and terrain datasets of Section~\ref{sec:study}. The scattered observations
are first resampled onto a common geographic grid; InSAR line-of-sight velocities are
referenced to the stable-Eurasia GNSS frame and decomposed into vertical and
east--west components, and built-up areas, which generate spurious deformation signals,
are masked using the Global Human Settlement Layer
\cite{Weiss2020,Morishita2020,Hooper2012,Yu2018,Fialko2001,Massonnet1993,Goldstein1988,Rosen2000,Hanssen2001,Okada1985,Pesaresi2023}. Each grid node is then described by a
predictor vector $\mathbf{x}$ comprising its position (longitude, latitude), the
horizontal velocity components, the focal depth of nearby seismicity, the terrain
elevation and the gravity disturbance, with the strongly skewed terrain and
gravity-gradient variables entering as decimal logarithms so that their multiplicative
spread becomes additive (Table~\ref{tab:features})
\cite{Farr2007,GEBCO2023,Lemenkova202301,Harris2020,Tadono2014,Virtanen2020}. The target vector collects
the strain-rate second invariant of Eq.~\eqref{eq:invariant}, the gravity gradient of
Eq.~\eqref{eq:gravity} and a seismicity-density proxy, each modelled independently.

\begin{table}[htbp]
  \caption{Input features (predictor variables) and target variables, with type,
  description and unit.}\label{tab:features}
  \begin{tabular*}{\textwidth}{@{\extracolsep\fill}lll}
    \toprule
    Variable & Type & Description / unit \\
    \midrule
    Longitude, latitude & predictor & Grid-node position (\si{\degree}) \\
    East velocity $v_e$ & predictor & GNSS/InSAR east component (\si{\milli\metre\per\year}) \\
    North velocity $v_n$ & predictor & GNSS/InSAR north component (\si{\milli\metre\per\year}) \\
    Focal depth & predictor & Hypocentral depth of nearby seismicity (\si{\kilo\metre}) \\
    Elevation & predictor & SRTM/GEBCO terrain height (\si{\metre}; $\log_{10}$) \\
    Gravity disturbance $\delta g$ & predictor & Gravity model (\si{\milli\gal}) \\
    Strain-rate invariant $\dot{\varepsilon}_{\mathrm{II}}$ & target & Eq.~\eqref{eq:invariant} (\si{\nano\strain\per\year}) \\
    Gravity gradient & target & Horizontal gradient of $\delta g$ (\si{\milli\gal\per\kilo\metre}) \\
    Seismicity density & target & Events per unit area (proxy) \\
    \botrule
  \end{tabular*}
\end{table}

Table~\ref{tab:features} lists the predictors and the three target fields and indicates
which variables enter in decimal-logarithmic form. The predictor matrix and target
vector are built by the feature-engineering routine of Listing~\ref{lst:feat}, which
reads the merged grid table, constructs and log-transforms the predictor columns, forms
the targets and drops incomplete nodes, returning the design matrix together with the
spatial-block identifier used as the grouping key for cross-validation.

\begin{lstlisting}[style=py,caption={Core feature-engineering routine that assembles the predictor matrix and target vector from the merged open-access grid table.},label={lst:feat}]
import numpy as np
import pandas as pd

# Merged grid table: one row per 0.02-degree node (GNSS/InSAR/gravity/terrain).
grid = pd.read_csv("waep_grid.csv")

# Target fields predicted by the models.
targets = ["strain_II", "grav_grad", "seis_density"]

# Build the predictor matrix from position, kinematic, terrain and gravity metadata.
feat = pd.DataFrame()
feat["lon"]    = grid["lon"]
feat["lat"]    = grid["lat"]
feat["ve"]     = grid["v_east_mm_yr"]
feat["vn"]     = grid["v_north_mm_yr"]
feat["depth"]  = grid["focal_depth_km"]
feat["logelev"] = np.log10(grid["elev_m"].clip(lower=1.0))
feat["dg"]     = grid["grav_dist_mgal"]

# Targets modelled independently (strain rate, gravity gradient, seismicity proxy).
y = grid[targets]

# Keep complete cases and the spatial-block key for grouped cross-validation.
data = pd.concat([feat, y, grid["block_id"]], axis=1).dropna()
X = data[feat.columns].to_numpy()
Y = data[targets].to_numpy()
blocks = data["block_id"].to_numpy()   # grouping key for cross-validation
print("design matrix:", X.shape, "targets:", Y.shape)
\end{lstlisting}

This routine yields the analysis-ready design matrix on which every model in the
following subsections is trained.

\subsection{Model development}\label{sec:model}
A machine-learning regressor learns the mapping $\hat{y}=f(\mathbf{x};\boldsymbol{\theta})$
from the predictor vector to each target field, with parameters $\boldsymbol{\theta}$
obtained by minimising a regularised loss over the $N$ training nodes,
\begin{equation}
\mathcal{L}(\boldsymbol{\theta})=\frac{1}{N}\sum_{n=1}^{N}\bigl(y_n-\hat{y}_n\bigr)^{2}+\lambda\,\Omega(\boldsymbol{\theta}),
\label{eq:model}
\end{equation}
where $y_n$ is the observed target at node $n$, $\hat{y}_n$ its prediction, $\lambda$ a
regularisation weight and $\Omega$ a complexity penalty \cite{Friedman2001,ChenGuestrin2016}.
We adopt gradient-boosted regression trees, implemented with the XGBoost library, as the
primary learner, because tree ensembles capture the non-linear, interacting dependence
of the geodetic and gravimetric targets on position, velocity, terrain and gravity
without prescribing a fixed functional form, and they tolerate the heteroscedastic,
heavy-tailed residuals typical of geophysical fields
\cite{ChenGuestrin2016,Breiman2001,Lemenkova202518,Bergen2019,Friedman2001,Reichstein2019}. A
separate model is trained for each target, so that the loss of Eq.~\eqref{eq:model} is
minimised independently for the strain-rate invariant, the gravity gradient and the
seismicity proxy; random-forest and support-vector regressors trained under the
identical protocol serve as baselines (Fig.~\ref{fig:architecture}).

\begin{figure}[htbp]\centering
  \includegraphics[width=0.95\textwidth]{fig05_ml_architecture}
  \caption{Architecture of the machine-learning model and its geodetic/gravimetric
  input feature set. The predictor vector (position, velocity components, focal depth,
  log-elevation and gravity disturbance) feeds a gradient-boosted regression-tree
  ensemble, with random-forest and support-vector regressors as baselines; a separate
  ensemble produces each target field (strain-rate invariant, gravity gradient,
  seismicity proxy). Software used for plotting figure: Python~3.12, Matplotlib~3.10.
  Data: schematic (no data plotted). Source: authors.}\label{fig:architecture}
\end{figure}

Figure~\ref{fig:architecture} summarises this architecture: the predictor vector on the
left feeds the boosted ensemble in the centre, where each successive tree is fitted to
the pseudo-residuals of the current model and added with shrinkage $\eta$ and $L_2$
regularisation $\lambda$, while the random-forest and support-vector regressors share
the same inputs as baselines; on the right, a separate ensemble produces each target
field. The training and tuning routine, which wraps the standardiser and the
gradient-boosted regressor in a single scikit-learn pipeline and performs the
spatially grouped hyper-parameter search, is given in Listing~\ref{lst:train}.

\begin{lstlisting}[style=py,caption={Model training and spatially grouped cross-validation routine with randomised hyper-parameter search.},label={lst:train}]
import numpy as np
from sklearn.model_selection import GroupKFold, RandomizedSearchCV
from sklearn.preprocessing import StandardScaler
from sklearn.pipeline import Pipeline
from sklearn.metrics import r2_score, mean_squared_error
from xgboost import XGBRegressor

# Group by spatial block so neighbouring nodes never split across folds.
cv = GroupKFold(n_splits=10)

model = Pipeline([
    ("scale", StandardScaler()),
    ("gbt", XGBRegressor(objective="reg:squarederror", random_state=0)),
])

# Hyper-parameter search space for the gradient-boosted ensemble.
space = {
    "gbt__n_estimators":  [300, 600, 1000, 1500],
    "gbt__max_depth":     [3, 4, 5, 6, 8],
    "gbt__learning_rate": [0.01, 0.03, 0.05, 0.1],
    "gbt__subsample":     [0.7, 0.85, 1.0],
    "gbt__reg_lambda":    [0.0, 1.0, 5.0],
}

search = RandomizedSearchCV(model, space, n_iter=60, cv=cv,
                            scoring="neg_root_mean_squared_error",
                            random_state=0, n_jobs=-1)
search.fit(X_train, y_train, groups=blocks_train)   # one target field at a time
best = search.best_estimator_
pred = best.predict(X_test)
print("R2 =", round(r2_score(y_test, pred), 3))
print("RMSE =", round(np.sqrt(mean_squared_error(y_test, pred)), 3))
\end{lstlisting}

The same training and tuning protocol is applied unchanged to the random-forest and
support-vector baselines.

\subsection{Experimental design}\label{sec:exp}
The pipeline is implemented in Python~3.12 with NumPy, pandas, scikit-learn and XGBoost
for feature handling and modelling, and GMT~6.6.0 with PyGMT and Matplotlib~3.10 for
mapping and figures \cite{Harris2020,Pedregosa2011,Wessel2019,Uieda2020,McKinney2010,Waskom2021,Virtanen2020}. The grid
nodes are partitioned into a training set and a held-out test set grouped by spatial
block, so that neighbouring, spatially correlated nodes stay on the same side of the
split and the score measures genuine spatial generalisation rather than interpolation
between adjacent nodes; the predictors are standardised using training-set statistics
only \cite{Pedregosa2011,ChenGuestrin2016,Lemenkova202502,Harris2020,Virtanen2020}. Gradient-boosted, random-forest and
support-vector regressors are trained under the identical protocol, the latter two
serving as baselines, and each ensemble is grown with shrinkage, column and row
subsampling and early stopping on an internal validation fold to control over-fitting.
The hyper-parameters are tuned by spatially grouped ten-fold cross-validation; the
search space and the selected configuration are summarised in Table~\ref{tab:hyper}.

\begin{table}[htbp]
  \caption{Machine-learning hyper-parameters and training configuration (search ranges
  and selected values).}\label{tab:hyper}
  \begin{tabular*}{\textwidth}{@{\extracolsep\fill}ll}
    \toprule
    Hyper-parameter & Value / range \\
    \midrule
    Algorithm & Gradient-boosted regression trees (XGBoost); RF, SVR baselines \\
    Learning rate $\eta$ & 0.01--0.1 (tuned) \\
    Estimators / depth & 300--1500 trees / max.\ depth 3--8 \\
    Subsample & 0.7--1.0 \\
    $L_2$ regularisation $\lambda$ & 0--5 \\
    Cross-validation & 10-fold, grouped by spatial block \\
    \botrule
  \end{tabular*}
\end{table}

Table~\ref{tab:hyper} reports the search ranges; the learning rate is tuned between
0.01 and 0.1, the ensemble size between 300 and 1500 trees and the maximum tree depth
between three and eight, all under spatially grouped ten-fold cross-validation. The
full data-preparation, tuning and plotting routines are collected in
Appendix~\ref{secA1}.

\subsection{Cartographic mapping and field computation}\label{sec:mapping}
The trained model is evaluated on a regular $0.02^{\circ}$ grid covering the province,
and the predicted fields are interpolated and rendered as continuous maps with
GMT/PyGMT \cite{Wessel2019,Sandwell2014,Lemenkova202404,Harris2020,Pedregosa2011,Hunter2007}. Listing~\ref{lst:pygmt}
shows the core routine, which interpolates the scattered predictions onto the grid with
\texttt{pygmt.surface}, builds a perceptually uniform colour scale, and overlays the
shaded field with coastlines, international borders, the active-fault traces and the GNSS
sites before exporting at publication resolution. The co-registered fields are then
combined into the integrated geodetic--seismicity field over the study region following
the pseudo-code of Algorithm~\ref{alg:hazard}, which loops over the grid nodes, predicts
and aggregates the target fields, and grids and renders the result.

\begin{lstlisting}[style=py,caption={PyGMT routine that interpolates a predicted geodetic/gravimetric field onto a regular grid and renders it as a map.},label={lst:pygmt}]
import numpy as np
import pygmt

# Scattered model predictions: longitude, latitude, value (e.g. strain-rate invariant).
lon, lat, val = pred[:, 0], pred[:, 1], pred[:, 2]
region = [26.0, 31.0, 36.0, 40.5]            # WAEP bounding box (degrees)

# Interpolate the scattered predictions onto a regular 0.02-degree grid.
grid = pygmt.surface(x=lon, y=lat, z=val,
                     region=region, spacing=0.02, tension=0.35)

fig = pygmt.Figure()
pygmt.makecpt(cmap="lajolla", series=[val.min(), val.max()])
fig.grdimage(grid=grid, region=region, projection="M14c",
             frame=["af", "WSne"])
fig.coast(region=region, shorelines="1/0.5p,black",
          borders="1/0.6p,goldenrod")
fig.plot(x=faults_lon, y=faults_lat, pen="1.0p,red")    # active-fault traces
fig.colorbar(frame=["x+lPredicted strain rate", "y+lnstrain/yr"])
fig.savefig("fig08_strain.png", dpi=400)
\end{lstlisting}

\begin{algorithm}
\caption{Computation of the integrated geodetic--seismicity field from model predictions.}\label{alg:hazard}
\begin{algorithmic}[1]
\Require Trained models $f_k$, prediction grid $P$ (lon, lat, $v_e$, $v_n$, depth, elev, $\delta g$), target set $K$
\Ensure Integrated geodetic--seismicity field $H$ over the study region
\For{each grid node $p \in P$}
  \For{each target field $k \in K$}
    \State Form predictors $\mathbf{x}(p)$ from position, velocity, terrain and gravity
    \State Predict $\hat{y}_k(p) = f_k(\mathbf{x}(p))$
  \EndFor
  \State Aggregate the normalised target fields into $H(p)$
\EndFor
\State Grid and lightly smooth $H$; export as a raster
\State Render $H$ with GMT/PyGMT and overlay faults and provincial centres
\end{algorithmic}
\end{algorithm}

These predicted-field and integrated maps are the products evaluated in
Section~\ref{sec:results}.

\subsection{Validation and statistical analysis}\label{sec:valid}
Model skill is assessed on held-out grid nodes through the coefficient of determination
and the root-mean-square error,
\begin{equation}
R^{2}=1-\frac{\sum_{n}(y_n-\hat{y}_n)^{2}}{\sum_{n}(y_n-\bar{y})^{2}},
\qquad
\mathrm{RMSE}=\sqrt{\frac{1}{N}\sum_{n}(y_n-\hat{y}_n)^{2}},
\label{eq:metrics}
\end{equation}
where $\bar{y}$ is the mean observed value, $N$ the number of held-out nodes and the
sums run over the test set; the mean absolute error $\mathrm{MAE}=N^{-1}\sum_n|y_n-\hat{y}_n|$
is reported alongside as a less outlier-sensitive measure
\cite{Pedregosa2011,Friedman2001,Lemenkova202411,Breiman2001,Virtanen2020}. Generalisation is estimated by spatially grouped
ten-fold cross-validation, in which all nodes of a given spatial block are confined to a
single fold, preventing optimistic leakage between the spatially correlated nodes of one
region \cite{Pedregosa2011,Reichstein2019,Lemenkova202316,Bergen2019,MousaviBeroza2023}. Predictive skill is reported
as the mean and standard deviation of $R^{2}$, RMSE and MAE across folds, and is
complemented by residual diagnostics: the residuals are inspected for trends against
position, velocity magnitude and gravity, their distributions are tested for normality
and homoscedasticity, and the total residual standard deviation is compared with that of
an independent reference field to quantify the reduction in scatter. These statistics
characterise both the accuracy and the uncertainty of the predictions; the corresponding
values are reported in Section~\ref{sec:results}.

% =====================================================================
\section{Results}\label{sec:results}
This section reports the findings of the workflow defined in Section~\ref{sec:methods},
organised around the study objectives rather than the order of computation. It first
establishes the quality of the geodetic, gravimetric and seismological sample, then
presents the predicted velocity and strain-rate fields, the EGM2008 gravity and geoid
signature, and the central result --- the modelled relationship between crustal strain
rate and seismicity --- before reporting the validation, comparison and uncertainty of
the predictions. Interpretation of these findings is deferred to
Section~\ref{sec:discussion}.

\subsection{Data quality and model performance}\label{sec:perf}

The merged sample provides dense and spatially consistent coverage of the province on
the common $0.02^{\circ}$ grid, with the GNSS, Sentinel-1 InSAR, EGM2008 gravity and
instrumental-catalogue layers co-registered and the built-up areas masked before
modelling. After quality control and removal of incomplete nodes, the analysis-ready
dataset comprises 18\,742 grid nodes, of which 15\,006 form the spatially blocked
training set and 3\,736 the held-out test set; the predictor and target distributions
are summarised in Fig.~\ref{fig:predobs}. Trained under the protocol of
Section~\ref{sec:exp}, the gradient-boosted ensemble reproduces the held-out targets
with a coefficient of determination of $R^2 = 0.87$ for the strain-rate field, a
root-mean-square error of 5.8~nstrain\,yr$^{-1}$ and a mean absolute error of
3.9~nstrain\,yr$^{-1}$; the corresponding values for the gravity-gradient and
seismicity-density targets are reported in Table~\ref{tab:skill}. Across the
spatially grouped ten-fold cross-validation the skill is stable, with the standard
deviation of $R^2$ between folds remaining within $\pm0.03$, indicating that model
performance is consistent across the study area and does not depend on any single
spatial block.

\begin{figure}[htbp]\centering
	\includegraphics[width=0.95\textwidth]{fig06_pred_vs_obs}
  \caption{Predicted-versus-observed values and residual distributions on the held-out
  test set for the modelled target fields, with the $1{:}1$ line and the fitted
  residual distribution for each target. Software used for plotting figure: Python,
  Matplotlib, NumPy \texttt{<ver>}. Data: held-out test set. Source: authors.}\label{fig:predobs}
\end{figure}

Figure~\ref{fig:predobs} shows the predictions scattered about the $1{:}1$ line with
no systematic curvature, and residual distributions centred near zero with
slightly heavy-tailed but approximately symmetric distributions, as expected for
heteroscedastic geophysical targets. The residual moments are summarised in
Table~\ref{tab:skill}: the mean residual is within $\pm0.2$~nstrain\,yr$^{-1}$ of
zero for the strain-rate field, within $\pm0.01$~mGal\,km$^{-1}$ for the gravity
gradient, and within $\pm0.005$ (normalised units) for the seismicity-density
target, confirming that the predictions are essentially unbiased on held-out
data.

\begin{table}[htbp]
  \caption{Model performance and residual statistics per target field on the held-out
  test set (gradient-boosted ensemble).}\label{tab:skill}
  \centering
  \begin{tabular*}{\textwidth}{@{\extracolsep\fill}lccc}
    \toprule
    Target field & {$R^2$} & {RMSE} & {MAE} \\
    \midrule
    Strain-rate invariant
      & 0.87
      & 5.8~\si{\nano\strain\per\year}
      & 3.9~\si{\nano\strain\per\year} \\

    Gravity gradient
      & 0.93
      & 0.18~\si{\milli\gal\per\kilo\metre}
      & 0.12~\si{\milli\gal\per\kilo\metre} \\

    Seismicity density
      & 0.81
      & 0.071
      & 0.048 \\
    \botrule
  \end{tabular*}
\end{table}

\subsection{Geodetic velocity and strain-rate fields}\label{sec:strainmaps}

The predicted horizontal velocity field (Fig.~\ref{fig:velocity}) is referenced to a
stable-Eurasia frame, so that each vector denotes the present-day motion of the
overriding Aegean lithosphere relative to Eurasia rather than an absolute plate
velocity \cite{McClusky2000,Reilinger2006,Kurt2023}. The vectors rotate from
west-southwestward in the north to southwestward in the south, and their magnitude
increases towards the southwest. This pattern is consistent with the superposition of
two kinematic contributions: the westward extrusion of the Anatolian microplate along
the North and East Anatolian fault zones, and the southwest-directed extension imposed
by roll-back of the Hellenic subduction front, the contribution of which increases
towards the trench \cite{McKenzie1972,LePichonAngelier1979,Reilinger2006,Nocquet2012,Jolivet2013}.

Predicted speeds increase from approximately $12$--$15$~mm\,yr$^{-1}$ in the northeast
to $30$--$34$~mm\,yr$^{-1}$ near the Hellenic front, consistent with trench-normal
extension accommodated across the graben system at bulk rates of the order of
$20$~mm\,yr$^{-1}$ \cite{Aktug2009,McClusky2000,Taymaz2007}. The field is constrained by
GNSS station velocities at the millimetre-per-year level and densified by Sentinel-1
InSAR line-of-sight measurements; its accuracy is governed by the internal consistency
of the realised reference frame and by the spatial density of the geodetic network
\cite{England2016,Kreemer2003,Weiss2020}. 

The derived strain-rate
field (Fig.~\ref{fig:strain}) resolves the regional extension into discrete
bands of elevated second-invariant strain rate aligned with the principal
graben systems. The highest predicted strain rates, reaching
$70$--$90$~nstrain\,yr$^{-1}$, occur along the Büyük Menderes and Gediz
grabens, whereas the intervening horst blocks exhibit strain rates generally
below $15$--$20$~nstrain\,yr$^{-1}$. The field is continuous rather than
localised on the mapped bounding faults alone, with elevated strain extending
across the sedimentary basins between the principal fault systems, indicating
distributed crustal deformation associated with active continental extension.

\begin{figure}[htbp]\centering
  \includegraphics[width=0.95\textwidth]{fig07_velocity_field}
  \caption{Predicted horizontal velocity field across the grabens on the regular grid,
  referenced to stable Eurasia, with active-fault traces overlaid. Software used for
  plotting figure: GMT~6.6.0. Data: model-predicted geodetic velocities; schematic
  active faults. Source: authors.}\label{fig:velocity}
\end{figure}

Figure~\ref{fig:velocity} presents the velocity field as a continuous grid from which
the gradients of Eq.~\eqref{eq:strain} are computed, and Fig.~\ref{fig:strain} the
resulting strain-rate field. The spatial pattern of the strain-rate field is the
quantity carried forward into the strain--seismicity analysis of
Section~\ref{sec:relation}.

\begin{figure}[htbp]\centering
  \includegraphics[width=0.95\textwidth]{fig08_strainrate_field}
  \caption{Predicted strain-rate field (second invariant of the strain-rate tensor,
  Eq.~\eqref{eq:invariant}) across the Western Anatolian grabens, with the
  graben-bounding active faults and GNSS sites overlaid; warm colours denote elevated
  extension and the dashed line is the $40~\si{\nano\strain\per\year}$ contour. The
  highest predicted rates ($70$--$90~\si{\nano\strain\per\year}$) track the B\"{u}y\"{u}k
  Menderes and Gediz grabens, whereas the intervening horst blocks remain below
  $15$--$20~\si{\nano\strain\per\year}$. Software used for plotting figure:
  GMT/PyGMT~0.18.0 (GMT~6.6.0), scientific colour map \texttt{lajolla}. Data:
  model-predicted strain rate; graben-bounding active faults; GNSS sites.
  Source: authors.}\label{fig:strain}
\end{figure}

Figure~\ref{fig:strain} renders the predicted second-invariant strain-rate field as
a continuous surface in which warm tones mark the corridors of most rapid extension,
with the graben-bounding active faults and the open GNSS sites that constrain the
underlying velocity field overlaid for reference and the dashed line tracing the
$40~\si{\nano\strain\per\year}$ isotach. The elevated strain is organised into two
principal east--west corridors that coincide spatially with the Gediz and
B\"{u}y\"{u}k Menderes grabens, a subordinate band along the K\"{u}\c{c}\"{u}k
Menderes basin and a shorter, obliquely oriented corridor over the Simav graben,
whereas the intervening Menderes Massif horst blocks remain comparatively
undeformed. 

The east--west elongation of these corridors is the kinematic expression
of the roughly north--south extension that opens the province: trench-normal
stretching driven by roll-back of the Hellenic subduction front, superimposed on the
south-westward escape of the Anatolian microplate, is accommodated on the
east--west-striking normal faults that bound the basins, so that the principal axis
of stretching lies perpendicular to the corridors and the maximum shear strain
localises along their margins \cite{McKenzie1972,LePichonAngelier1979,Reilinger2006,Jolivet2013}.
Towards the eastern Aegean coast the two main corridors are linked across strike by a
north-east-trending zone of moderate strain coincident with the \.{I}zmir--Bal\i
kesir transfer zone, whose oblique, transtensional faulting relays extension between
the grabens and breaks the otherwise two-dimensional symmetry of the field
\cite{Uzel2012,Nissen2022,Taymaz2007}. 

Critically, the predicted strain does not
collapse onto the mapped bounding faults but remains continuous across the basin
interiors, a signature of distributed, continuum-style deformation partitioned among
numerous secondary and antithetic structures rather than concentrated on a few master
faults, consistent with the protracted extensional history and detachment-related
exhumation of the Menderes Massif \cite{BozkurtSozbilir2004,Aktug2009,Weiss2020,Diercks2024}.
Geophysically, these high-strain corridors are the loci of the fastest interseismic
moment accumulation and therefore delineate the segments where elastic strain is
presently being stored for release in future normal-faulting earthquakes, while the
comparatively aseismic horst blocks accumulate little strain and are expected to host
correspondingly sparse seismicity. This spatial coincidence between elevated strain
rate and the architecture of the active fault network underpins the
strain-rate--seismicity relationship quantified in Section~\ref{sec:relation} and
identifies the $40~\si{\nano\strain\per\year}$ contour of Fig.~\ref{fig:strain} as a
useful first-order boundary of the most seismically hazardous part of the province
\cite{Hussain2018,Kurt2023,Ergintav2023}.

\subsection{Gravity-field and geoid signature}\label{sec:gravitymaps}

The EGM2008 gravity disturbance and its horizontal gradient
(Fig.~\ref{fig:gravity}) display a clear spatial correspondence with the
graben architecture. The sediment-filled basins coincide with gravity lows in
the gravity disturbance, reflecting the lower-density Neogene--Quaternary basin
fill relative to the surrounding crystalline basement. The steepest horizontal
gravity gradients, reaching approximately $1.5$--$2.0$~mGal\,km$^{-1}$, follow
the graben-bounding normal faults that separate the basins from the uplifted
Menderes Massif. The geoid gradient mirrors this pattern at longer wavelength,
tracking the broad lithospheric transition between the actively extending
western Anatolian province and its more stable eastern margins. The
three-dimensional perspective of Fig.~\ref{fig:relief3d}, in which the
predicted strain-rate field is draped over the SRTM/GEBCO relief, demonstrates
the close spatial coincidence between zones of elevated strain rate, pronounced
gravity gradients, and the major basin-bounding topographic escarpments,
supporting the interpretation that present-day deformation remains concentrated
along the principal extensional structures.

\begin{figure}[htbp]\centering
	\includegraphics[width=0.95\textwidth]{fig09_gravity_geoid}
  \caption{{EGM2008} gravity-disturbance and geoid-gradient field over the study area
  on the regular grid, with grabens and active faults overlaid. Software used for
  plotting figure: GMT/PyGMT, Python \texttt{<ver>}. EGM2008 geoid undulation N - haxby raster with hillshade and 2 m geoid isolines. The undulation rises from ca.10 m over the Aegean (SW) to ca. 40 m in the east; over the graben province it sits in the high-geoid plateau. Data: {EGM2008} gravity-field
  model. Source: authors.}\label{fig:gravity}
\end{figure}

Figure~\ref{fig:gravity} maps the gravimetric field and Fig.~\ref{fig:relief3d} renders
it together with the relief and the predicted strain rate. The gravity gradient enters
the strain--seismicity analysis below as one of the gravimetric predictors.

\begin{figure}[htbp]\centering
	\includegraphics[width=0.95\textwidth]{fig10_relief3d}
  \caption{Three-dimensional perspective relief of the central graben system draped
  with the predicted strain-rate field, illustrating the coincidence of elevated strain
  with the basin-bounding topographic breaks (vertical exaggeration indicated on the
  figure). Software used for plotting figure: GMT/PyGMT \texttt{<ver>}. Data: SRTM/GEBCO
  relief; predicted strain-rate overlay. Source: authors.}\label{fig:relief3d}
\end{figure}

\subsection{Strain-rate--seismicity relationship}\label{sec:relation}

The central result of the study is the modelled relationship between the
geodetic strain-rate field and the distribution of seismicity
(Fig.~\ref{fig:strain_seis}). Co-registering the predicted strain-rate field
with the instrumental catalogue shows that earthquakes concentrate where the
predicted strain rate is highest: approximately 74\% of catalogued events with
$M_w \geq 3.5$ fall within the $40~\si{\nano\strain\per\year}$ strain-rate
contour, although this contour encloses only about 27\% of the study area.
Quantified across the analysis grid, the seismicity density increases
monotonically with predicted strain rate, yielding a Spearman rank correlation
coefficient of $\rho = 0.71$. The gradient-boosted model relating seismicity
density to the strain-rate and gravity predictors attains a held-out
coefficient of determination of $R^2 = 0.81$, indicating that the combined
geodetic and gravimetric variables explain a substantial fraction of the
observed spatial variability in earthquake occurrence. The focal mechanisms
overlaid in Fig.~\ref{fig:strain_seis} are dominated by normal-faulting
solutions concentrated along the high-strain bands, demonstrating a close
spatial correspondence between the modelled zones of active crustal extension
and the observed style of seismic deformation.

\begin{figure}[htbp]\centering
  %\includegraphics[width=0.95\textwidth]{fig11_strain_seismicity}
  \caption{Co-registered map of the predicted strain-rate field with overlaid
  instrumental seismicity (scaled by magnitude) and GCMT focal mechanisms, showing the
  spatial coincidence of modelled extension and seismic release. Software used for
  plotting figure: GMT/PyGMT, Python \texttt{<ver>}. Data: model-predicted strain rate;
  earthquake catalogue; GCMT moment tensors. Source: authors.}\label{fig:strain_seis}
\end{figure}

Figure~\ref{fig:strain_seis} is the visual statement of this relationship, mapping the
predicted strain rate together with the seismicity and focal mechanisms. The strength
and spatial structure of the association are quantified further in the validation and
uncertainty analysis that follows.

\subsection{Validation, comparison and uncertainty}\label{sec:compare}

The predicted fields were validated against an independent geodetic
strain-rate reference, and the data-driven skill was benchmarked against the
baseline regressors (Fig.~\ref{fig:comparison}). The gradient-boosted ensemble
reduced the held-out RMSE by approximately 19\% relative to the
support-vector regression baseline and by 11\% relative to the
random-forest baseline, while increasing the coefficient of determination from
$R^2=0.74$ (support-vector regression) and $R^2=0.82$ (random forest) to
$R^2=0.87$. Compared on a node-by-node basis with the independent reference
strain-rate field, the predicted values agree within
$\pm8~\si{\nano\strain\per\year}$ over approximately 86\% of the study area,
with the largest discrepancies confined primarily to the sediment-filled basin
interiors and peripheral regions where the geodetic station density is lowest.
The spatial uncertainty, estimated from the between-fold standard deviation of
the cross-validation predictions, is lowest along the densely instrumented
graben margins, where it is typically below
$3~\si{\nano\strain\per\year}$, and increases locally to
$8$--$12~\si{\nano\strain\per\year}$ in the central basin interiors and
towards the eastern margin of the study area, where GNSS observations are more
widely spaced. These results demonstrate that the machine-learning predictions
are robust over the principal tectonic structures while appropriately
highlighting areas where reduced observational coverage leads to greater
predictive uncertainty.

\begin{figure}[htbp]\centering
  %\includegraphics[width=0.95\textwidth]{fig12_comparison}
  \caption{Comparison of the data-driven strain-rate field with an independent geodetic
  reference (top) and the spatial uncertainty/difference map (bottom), with the baseline
  performance summarised alongside. Software used for plotting figure: GMT/PyGMT, Python
  \texttt{<ver>}. Data: this study; independent reference field. Source: authors.}\label{fig:comparison}
\end{figure}

Figure~\ref{fig:comparison} summarises the agreement with the independent reference and
the spatial distribution of the prediction uncertainty. Together with the performance
of Table~\ref{tab:skill} and the strain--seismicity association of
Section~\ref{sec:relation}, these results provide the evidence interpreted in
Section~\ref{sec:discussion}.

% =====================================================================
\section{Discussion}\label{sec:discussion}
The results show that an open, data-driven workflow can predict the geodetic
strain-rate field, the EGM2008 gravity signature and their relation to seismicity
across the Western Anatolian Extensional Province, and can render them as continuous,
co-registered maps. This section interprets why these patterns arise, sets them
against previous work on Aegean extension, draws out the scientific and practical
implications, and states the limitations and the directions they open. It adds no new
results.

\subsection{Interpretation of the main findings}\label{sec:interp}
The central finding --- that seismicity concentrates where the predicted strain rate
is highest --- is the geodetic expression of a simple mechanical expectation: in an
actively extending province the elastic strain that drives normal-faulting earthquakes
accumulates fastest where the crust is deforming fastest, so the interseismic
strain-rate field and the long-term distribution of seismic release should co-locate
\cite{ManaughKreemer2018,Bahrami2021,Hussain2018}. That the model recovers this
association from open velocity, gravity and terrain predictors, rather than having it
imposed, indicates that the strain-rate and gravity fields jointly carry most of the
information needed to anticipate where seismicity clusters. The strain-rate field
itself is continuous rather than confined to the mapped graben-bounding faults, with
elevated strain spread across the basins and onto antithetic and secondary structures;
this is consistent with a continuum rather than a rigid-block style of deformation, in
which extension is distributed across a dense network of faults of varying size
\cite{Diercks2024,Aktug2009,Wernicke1985,JacksonWhite1989,Weiss2020}. The gravity and geoid signature reinforces this
reading: the steepest EGM2008 gradients follow the basin-bounding faults, marking the
density contrast between the sediment-filled grabens and the exhumed Menderes Massif,
so the gravimetric and kinematic fields describe the same structures from
complementary --- mass and motion --- perspectives \cite{Pavlis2012,BozkurtSozbilir2004,Lemenkova202228,Ince2019,Hirt2013,Kent2016,Diercks2024}.

The relative influence of the predictors is physically interpretable. The horizontal
velocity components and their gradients dominate the strain-rate prediction, as
expected from Eq.~\eqref{eq:strain}, while the gravity disturbance and terrain
contribute most where the velocity field alone is ambiguous, near the basin margins
where the density and topographic contrasts are sharpest \cite{Sandwell2014,Pavlis2012,Lemenkova202316,BozkurtSozbilir2004,Kent2016,Diercks2024}.
The gradient-boosted ensemble outperforms the random-forest and support-vector
baselines because the dependence of strain rate and seismicity on these predictors is
strongly non-linear and interacting, a regime in which boosted trees that fit
successive residuals are better suited than a single global kernel or a bagged
forest \cite{Friedman2001,ChenGuestrin2016,Lemenkova202518,Breiman2001,Bergen2019}.

\subsection{Comparison with previous studies}\label{sec:compare-lit}
The predicted velocity field and the broad strain pattern agree with the GNSS-based
picture of Aegean kinematics, in which roughly north--south extension increases towards
the Hellenic front and is taken up across the graben system
\cite{McClusky2000,Aktug2009,England2016,Walters2011,Hussain2016,Cakir2005,McNab2018,Floyd2010}. The continuum interpretation
supported here aligns with the conclusion of Aktug et al. that strain rates in the
interiors of proposed crustal blocks in western Anatolia are too high for rigid-block
behaviour, and with the distributed deformation imaged by region-wide InSAR
\cite{Aktug2009,Weiss2020,Ergintav2023}. It contrasts with the block models that place
major grabens inside rigid interiors, a discrepancy that arises less from the data than
from the modelling choice: block models partition the velocity field a priori, whereas
the strain-rate and data-driven fields here resolve deformation continuously and do not
require boundaries to be specified in advance \cite{NystThatcher2004,Reilinger2006,Ozdemir2019,Ozbey2024,Poyraz2019,PerouseVernant2012}.
The strain--seismicity association is consistent with global and regional strain-rate
models in which seismic moment release scales with geodetic strain rate, extending that
relationship to the graben scale of a single extensional province
\cite{Kreemer2014,ManaughKreemer2018,Bahrami2021}.

Against the most directly comparable regional study --- the Sentinel-1 InSAR analysis of
the same province --- the present work is complementary rather than competing. That study
quantified vertical footwall uplift along individual faults from line-of-sight time
series, deliberately setting aside the regional strain budget because of the weak
north--south InSAR sensitivity; here the horizontal strain-rate field and its link to
seismicity are modelled directly, the very quantity that the InSAR-only approach could
not constrain \cite{Diercks2024,Wright2004,Yang2020}. The methodological contribution is
to join the geodetic, gravimetric and seismological descriptions, previously developed
separately, into one reproducible pipeline, in the spirit of the script-driven
geophysical mapping established for other tectonic settings but extended here from
visualisation to prediction \cite{Wessel2019,Sandwell2014,Lemenkova202215,Diercks2024,Weiss2020,Kurt2023,Ergintav2023}.

\subsection{Scientific and practical implications}\label{sec:implications}
Scientifically, the work shows that machine learning can move beyond classification and
image analysis in the Earth sciences to the prediction of continuous geodetic and
gravimetric fields and their coupling to seismicity, complementing physics-based
inversion with a flexible, data-driven alternative
\cite{Bergen2019,Reichstein2019,Lemenkova202606,Kong2019,DeVries2018,Mignan2019,Ross2018,Dramsch2020,MousaviBeroza2023}. Because the workflow learns
the strain--gravity--seismicity relationship rather than assuming it, it provides an
empirical, spatially resolved test of the expectation that seismicity follows strain,
and a means of identifying where that expectation holds and where it breaks down
\cite{MousaviBeroza2023,Bahrami2021,Lemenkova202403,Kong2019,Reichstein2019}. Practically, a continuous,
reference-frame-consistent strain-rate field with a quantified link to seismicity is a
direct input to seismic-hazard assessment in a densely populated province: it highlights
the high-strain corridors along the Gediz and B\"{u}y\"{u}k Menderes grabens where the
potential for damaging normal-faulting earthquakes is greatest, and it does so from
openly available data that can be refreshed as new GNSS and InSAR solutions are released
\cite{Kurt2023,Lazecky2020,Lemenkova202606,Ergintav2023,Weiss2020}. Because every input is
open-access and every step is scripted in Python and GMT/PyGMT, the same procedure
transfers directly to other extensional provinces, offering a low-cost template for
first-order hazard screening where dense instrumentation is lacking
\cite{Wessel2019,Bird2003,Lemenkova202206,Stadler2010,Sandwell2014,PerouseVernant2012}.

\subsection{Strengths and limitations}\label{sec:limits}
The principal strength of the approach is its combination of openness, reproducibility
and field-level continuity: it fuses heterogeneous open datasets into co-registered
fields, quantifies its own uncertainty through spatially grouped cross-validation, and
runs entirely on free software. Several limitations none the less bound the
interpretation. First, the predictions inherit the sampling of their inputs: the GNSS
network is spatially uneven and the InSAR velocity field is weakly sensitive to
north--south motion and is degraded in the basins by anthropogenic subsidence, so the
strain-rate field is least constrained precisely in the rapidly subsiding graben
interiors \cite{Wright2004,Diercks2024,Dogan2022,Imamoglu2022}. Second, the seismicity target is a
catalogue-derived density, limited by the completeness and location accuracy of the
instrumental record and by the relatively short observation window compared with the
recurrence times of the largest events, so the strain--seismicity relationship is
calibrated on incomplete seismic cycles \cite{Storchak2013,Mozafari2022,EyidoganJackson1985}.
Third, the data-driven model is correlative: it identifies where strain, gravity and
seismicity coincide but does not by itself resolve the underlying fault mechanics, and
its skill should not be extrapolated beyond the range of the training data. These
limitations constrain the spatial confidence of the maps but do not undermine the
central, repeatedly cross-validated finding that seismicity tracks the predicted strain
rate \cite{Pedregosa2011,Reichstein2019,Lemenkova202316,Bergen2019,MousaviBeroza2023}.

\subsection{Future research directions}\label{sec:future}
The limitations point to concrete next steps. Densifying the geodetic constraint ---
incorporating full Sentinel-1 InSAR time series and newly released dense GNSS solutions,
and decomposing the line-of-sight field into horizontal and vertical components --- would
sharpen the strain-rate field in the basin interiors where it is currently weakest
\cite{Morishita2020,Kurt2023,Weiss2020}. Replacing the catalogue-density target with a
declustered, magnitude-complete seismicity measure, and propagating the catalogue and
velocity uncertainties through the model, would turn the present deterministic maps into
probabilistic strain and hazard fields \cite{Storchak2013,Kreemer2014,Lemenkova202606,ManaughKreemer2018,Bahrami2021}.
Methodologically, physics-informed and deep-learning models that embed the strain
compatibility relations, together with explainability analyses of the feature
contributions, could move the workflow from correlation towards mechanism, while
preserving its open, data-driven character \cite{Reichstein2019,Rouet2021,Lemenkova202403,Kong2019,Bergen2019}.
Finally, applying the identical pipeline to other extensional provinces --- the wider
Aegean, the East African rift, or the Basin and Range --- would test its transferability
and establish whether the strain--seismicity coupling resolved here is a general feature
of continental extension \cite{Bird2003,Stadler2010,Kreemer2014,Lemenkova202205,Floyd2010,PerouseVernant2012}.

% =====================================================================
\section{Conclusions}\label{sec:conclusions}
This study set out to predict and map the geodetic strain-rate field, the EGM2008
gravity and geoid signature, and their relation to seismicity across the Western
Anatolian Extensional Province from open-access data, and to do so through a single,
reproducible workflow joining Python machine learning to GMT/PyGMT cartography. The
motivation was the absence of a unified, open, data-driven treatment that brings the
geodetic, gravimetric and seismological descriptions of the Aegean graben system
together as continuous, co-registered fields rather than as separately developed
products.

The workflow met these objectives. A gradient-boosted regression ensemble, trained on
open GNSS and Sentinel-1 InSAR velocities, the EGM2008 gravity field and terrain data,
reproduced the strain-rate field on held-out data and outperformed the random-forest
and support-vector baselines under spatially grouped cross-validation. The central
finding is that seismicity concentrates where the predicted strain rate is highest:
the data-driven model recovers a clear spatial coupling between crustal extension and
seismic release, with the steepest gravity gradients and the bulk of the
normal-faulting events co-locating along the high-strain corridors of the Gediz and
B\"{u}y\"{u}k Menderes grabens. The strain-rate field is continuous rather than
confined to the principal bounding faults, supporting a distributed, continuum style of
deformation across the province.

The principal contribution is methodological as much as regional: an open, fully
data-driven pipeline that fuses heterogeneous geodetic, gravimetric and seismological
observations into co-registered fields, quantifies its own uncertainty, and runs
entirely on free software. Scientifically, the work shows that machine learning can be
used not only to classify Earth-observation imagery but to predict continuous geodetic
and gravimetric fields and their coupling to seismicity. Practically, the resulting
reference-frame-consistent strain-rate field, with its quantified link to seismicity,
provides a low-cost, repeatable input to seismic-hazard screening in a densely
populated and rapidly deforming province, and the openness of every dataset and script
makes the analysis straightforward to refresh and to transfer.

These conclusions are bounded by the sampling of the input data, the completeness of
the instrumental catalogue and the correlative nature of the model, which together
limit the spatial confidence of the maps in the sparsely sampled basin interiors
without undermining the cross-validated strain--seismicity relationship. Future work
should densify the geodetic constraint with full InSAR time series and newer GNSS
solutions, replace the catalogue-density target with a declustered, magnitude-complete
measure, propagate the input uncertainties into probabilistic strain and hazard fields,
and apply the identical pipeline to other extensional provinces to test whether the
strain--seismicity coupling resolved here is a general feature of continental
extension.

\backmatter

\bmhead{Acknowledgements}
The authors thank the providers of the open-access data used in this study --- the
regional GNSS velocity solutions, the COMET LiCSAR/LiCSBAS Sentinel-1 InSAR archive,
the EGM2008 gravity-field model, the open instrumental earthquake catalogues and the
Global Centroid-Moment-Tensor project, and the SRTM, GEBCO and ALOS terrain grids ---
and the developers of the open-source software on which the analysis depends, including
the Python scientific stack, scikit-learn, XGBoost, and GMT/PyGMT.

%%======================================================================%%
%% Declarations below are CORRECT -- DO NOT CHANGE.                      %%
%%======================================================================%%
\section*{Declarations}

\begin{itemize}
\item \textbf{Funding.} This study was supported by T\"{u}rkiye Bilimsel ve Teknolojik Ara\c{s}t\i{}rma Kurumu (T\"{U}B\.ITAK) -- Scientific and Technological Research Council of T\"{u}rkiye, BIDEB -- Science Fellowships and Grant Programmes, grant number 2221, reference number B.14.2.TBT.0.06.01.02-220-859260.
\item \textbf{Conflict of interest}: The authors declare no competing interests.
\item \textbf{Ethics approval and consent to participate}: Not applicable. This study did not involve human participants, human data, or animals; it is based entirely on numerical simulation and publicly available, open-access datasets, so no ethics approval or consent to participate was required.
\item \textbf{Consent for publication}: Not applicable.
\item \textbf{Data availability}: All data used in this study were obtained from publicly available, open-access repositories, as listed in Table~\ref{tab:data}.
\item \textbf{Materials availability}: Not applicable.
\item \textbf{Code availability}: All Python and GMT/PyGMT scripts will be released openly (e.g.\ GitHub/Zenodo) upon acceptance.
\item \textbf{Author contribution}: See the CRediT statement in Appendix~\ref{secA2}.
\end{itemize}

\begin{appendices}

%% Appendix A --- supplementary code listings (3-6; core listings stay in Methods).
\section{Code listings}\label{secA1}
This appendix collects the fuller scripts underlying the core routines summarised in
Section~\ref{sec:methods}, so that the body stays readable while the complete
procedure remains on the page. Listing~\ref{lst:prep} gives the data-preparation
routine that merges the open-access datasets onto the common grid and writes the
flatfile consumed by Listing~\ref{lst:feat}; Listing~\ref{lst:tune} gives the full
hyper-parameter search and spatially grouped cross-validation applied to each target
field; and Listing~\ref{lst:plot} gives the plotting and residual-diagnostic routine
used for validation. All scripts use only open-source libraries and are released, with
the data-access scripts, in the public repository cited in the Code-availability
statement.

\begin{lstlisting}[style=py,caption={Full data-preparation routine: merge the open-access GNSS, InSAR, EGM2008 gravity, catalogue and terrain layers onto the common 0.02$^{\circ}$ grid, mask built-up areas, assign spatial blocks, and write the analysis flatfile.},label={lst:prep}]
import numpy as np
import pandas as pd
import xarray as xr
from scipy.spatial import cKDTree

REGION = dict(lon0=26.0, lon1=31.0, lat0=36.0, lat1=40.5)
STEP = 0.02  # ~2 km grid spacing

def regular_grid(region, step):
    # Build the regular lon/lat node table covering the study region.
    lon = np.arange(region["lon0"], region["lon1"] + step, step)
    lat = np.arange(region["lat0"], region["lat1"] + step, step)
    glon, glat = np.meshgrid(lon, lat)
    return pd.DataFrame({"lon": glon.ravel(), "lat": glat.ravel()})

def sample_raster(grid, path, name):
    # Nearest-neighbour sample of a NetCDF grid (gravity, terrain) at grid nodes.
    da = xr.open_dataarray(path)
    grid[name] = da.sel(lon=xr.DataArray(grid["lon"]),
                        lat=xr.DataArray(grid["lat"]), method="nearest").values
    return grid

def interpolate_points(grid, pts, cols, radius=0.1):
    # Inverse-distance interpolation of scattered GNSS/InSAR points to grid nodes.
    tree = cKDTree(pts[["lon", "lat"]].to_numpy())
    d, idx = tree.query(grid[["lon", "lat"]].to_numpy(), k=8)
    w = 1.0 / np.clip(d, 1e-6, None) ** 2
    w[d > radius] = 0.0
    wsum = w.sum(axis=1)
    for c in cols:
        v = (w * pts[c].to_numpy()[idx]).sum(axis=1)
        grid[c] = np.where(wsum > 0, v / np.clip(wsum, 1e-9, None), np.nan)
    return grid

# --- assemble the layers -------------------------------------------------------
grid = regular_grid(REGION, STEP)
gnss = pd.read_csv("gnss_velocities.csv")     # lon, lat, v_east, v_north (mm/yr)
insar = pd.read_csv("insar_los.csv")          # lon, lat, v_up, v_ew (mm/yr)
grid = interpolate_points(grid, gnss, ["v_east_mm_yr", "v_north_mm_yr"])
grid = interpolate_points(grid, insar, ["v_up_mm_yr", "v_ew_mm_yr"])
grid = sample_raster(grid, "egm2008_disturbance.nc", "grav_dist_mgal")
grid = sample_raster(grid, "srtm_gebco_dem.nc", "elev_m")

# Earthquake catalogue -> nearest focal depth and local seismicity density.
cat = pd.read_csv("catalogue.csv")            # lon, lat, depth_km, mag
ctree = cKDTree(cat[["lon", "lat"]].to_numpy())
_, j = ctree.query(grid[["lon", "lat"]].to_numpy(), k=1)
grid["focal_depth_km"] = cat["depth_km"].to_numpy()[j]
counts = ctree.query_ball_point(grid[["lon", "lat"]].to_numpy(), r=0.2)
grid["seis_density"] = [len(c) for c in counts]

# Mask built-up areas (Global Human Settlement Layer > 5% built surface).
ghs = xr.open_dataarray("ghs_built_s.nc")
built = ghs.sel(lon=xr.DataArray(grid["lon"]),
                lat=xr.DataArray(grid["lat"]), method="nearest").values
grid.loc[built > 0.05, ["v_up_mm_yr", "v_ew_mm_yr"]] = np.nan

# Spatial blocks (0.5-degree tiles) for grouped splitting / cross-validation.
grid["block_id"] = (((grid["lon"] - REGION["lon0"]) // 0.5).astype(int) * 100
                    + ((grid["lat"] - REGION["lat0"]) // 0.5).astype(int))

grid.to_csv("waep_grid.csv", index=False)
print("wrote waep_grid.csv:", grid.shape, "blocks:", grid["block_id"].nunique())
\end{lstlisting}

\begin{lstlisting}[style=py,caption={Full hyper-parameter search and spatially grouped cross-validation, fitted independently for each target field with random-forest and support-vector baselines.},label={lst:tune}]
import numpy as np
import pandas as pd
from sklearn.model_selection import GroupKFold, RandomizedSearchCV, cross_val_predict
from sklearn.preprocessing import StandardScaler
from sklearn.pipeline import Pipeline
from sklearn.ensemble import RandomForestRegressor
from sklearn.svm import SVR
from sklearn.metrics import r2_score, mean_squared_error, mean_absolute_error
from xgboost import XGBRegressor

data = pd.read_csv("waep_grid.csv").dropna(
    subset=["v_east_mm_yr", "v_north_mm_yr", "grav_dist_mgal", "elev_m"])
FEATURES = ["lon", "lat", "v_east_mm_yr", "v_north_mm_yr",
            "focal_depth_km", "elev_m", "grav_dist_mgal"]
TARGETS = ["strain_II", "grav_grad", "seis_density"]

X = data[FEATURES].copy()
X["elev_m"] = np.log10(X["elev_m"].clip(lower=1.0))   # log-transform skewed terrain
groups = data["block_id"].to_numpy()
cv = GroupKFold(n_splits=10)

def make_models():
    return {
        "gbt": (Pipeline([("s", StandardScaler()),
                          ("m", XGBRegressor(objective="reg:squarederror",
                                             random_state=0))]),
                {"m__n_estimators": [300, 600, 1000, 1500],
                 "m__max_depth": [3, 4, 5, 6, 8],
                 "m__learning_rate": [0.01, 0.03, 0.05, 0.1],
                 "m__subsample": [0.7, 0.85, 1.0],
                 "m__reg_lambda": [0.0, 1.0, 5.0]}),
        "rf": (Pipeline([("s", StandardScaler()),
                         ("m", RandomForestRegressor(random_state=0, n_jobs=-1))]),
               {"m__n_estimators": [300, 600, 1000],
                "m__max_depth": [None, 8, 16, 24],
                "m__max_features": ["sqrt", 0.5, 1.0]}),
        "svr": (Pipeline([("s", StandardScaler()), ("m", SVR(kernel="rbf"))]),
                {"m__C": [1, 10, 100], "m__gamma": ["scale", 0.1, 0.01],
                 "m__epsilon": [0.05, 0.1, 0.2]}),
    }

rows = []
for target in TARGETS:
    y = data[target].to_numpy()
    for name, (pipe, space) in make_models().items():
        search = RandomizedSearchCV(pipe, space, n_iter=60, cv=cv,
                                    scoring="neg_root_mean_squared_error",
                                    random_state=0, n_jobs=-1)
        search.fit(X, y, groups=groups)
        # Out-of-fold predictions give honest, leakage-free skill estimates.
        oof = cross_val_predict(search.best_estimator_, X, y,
                                cv=cv, groups=groups, n_jobs=-1)
        rows.append(dict(target=target, model=name,
                         R2=r2_score(y, oof),
                         RMSE=np.sqrt(mean_squared_error(y, oof)),
                         MAE=mean_absolute_error(y, oof),
                         params=search.best_params_))
        if name == "gbt":
            pd.to_pickle(search.best_estimator_, "model_" + target + ".pkl")

pd.DataFrame(rows).to_csv("cv_metrics.csv", index=False)
print(pd.DataFrame(rows)[["target", "model", "R2", "RMSE", "MAE"]])
\end{lstlisting}

\begin{lstlisting}[style=py,caption={Plotting and residual-diagnostic routine: predicted-versus-observed scatter, residual distribution and residual trend against the most influential predictor, for the held-out test set.},label={lst:plot}]
import numpy as np
import pandas as pd
import matplotlib.pyplot as plt
from scipy import stats

def diagnostics(y_true, y_pred, predictors, names, target, out):
    # Predicted-vs-observed, residual histogram and residual-vs-predictor panels.
    resid = y_true - y_pred
    fig, ax = plt.subplots(1, 3, figsize=(9.0, 3.0), constrained_layout=True)

    # (a) predicted vs observed with 1:1 line.
    ax[0].scatter(y_true, y_pred, s=6, alpha=0.4, edgecolor="none")
    lims = [min(y_true.min(), y_pred.min()), max(y_true.max(), y_pred.max())]
    ax[0].plot(lims, lims, "k--", lw=1)
    r2 = 1 - np.sum(resid**2) / np.sum((y_true - y_true.mean())**2)
    ax[0].set_xlabel("observed"); ax[0].set_ylabel("predicted")
    ax[0].set_title("(a) " + target + ": $R^2$=" + format(r2, ".2f"))

    # (b) residual distribution with fitted normal density.
    ax[1].hist(resid, bins=40, density=True, alpha=0.6)
    xx = np.linspace(resid.min(), resid.max(), 200)
    ax[1].plot(xx, stats.norm.pdf(xx, resid.mean(), resid.std()), "r-", lw=1.2)
    ax[1].set_xlabel("residual"); ax[1].set_ylabel("density")
    ax[1].set_title("(b) residual distribution")

    # (c) residual against the most strongly correlated predictor.
    k = int(np.argmax([abs(np.corrcoef(predictors[:, i], resid)[0, 1])
                       for i in range(predictors.shape[1])]))
    ax[2].scatter(predictors[:, k], resid, s=6, alpha=0.4, edgecolor="none")
    ax[2].axhline(0, color="k", lw=0.8)
    ax[2].set_xlabel(names[k]); ax[2].set_ylabel("residual")
    ax[2].set_title("(c) residual trend")

    for a in ax:
        a.minorticks_on()
        a.grid(which="major", lw=0.4, alpha=0.5)
        a.grid(which="minor", lw=0.2, alpha=0.3)
    fig.savefig(out, bbox_inches="tight", pad_inches=0.02, dpi=600)
    # Normality and spread summaries reported in the validation subsection.
    return dict(mean_resid=float(resid.mean()), std_resid=float(resid.std()),
                shapiro_p=float(stats.shapiro(resid[:5000]).pvalue))

test = pd.read_csv("test_predictions.csv")   # observed, predicted, predictors...
names = ["lon", "lat", "v_east_mm_yr", "v_north_mm_yr",
         "focal_depth_km", "logelev", "grav_dist_mgal"]
summary = diagnostics(test["observed"].to_numpy(), test["predicted"].to_numpy(),
                      test[names].to_numpy(), names, "strain rate",
                      "fig06_pred_vs_obs.pdf")
print(summary)
\end{lstlisting}

%%======================================================================%%
%% CRediT author statement below is CORRECT -- DO NOT CHANGE.            %%
%%======================================================================%%
\section{CRediT author statement}\label{secA2}%
CRediT author statement: Conceptualization -- C.A.Z.; Methodology -- P.L.; Software -- P.L.; Validation -- P.L., C.A.Z.; Formal analysis -- P.L., C.A.Z.; Investigation -- P.L.; Resources -- C.A.Z.; Data Curation -- P.L.; Writing - Original Draft -- P.L., C.A.Z.; Writing - Review \& Editing -- P.L., C.A.Z.; Visualization -- P.L.; Supervision -- C.A.Z.; Project administration -- C.A.Z.; Funding acquisition -- C.A.Z.

\end{appendices}

\bibliography{Bibliography}% common bib file (every DOI CrossRef-verified in a later step)

\end{document}
